%! Graphing Rational Functions — Unit 5, Lesson 5.5 — Exit Ticket (Answer Key)
\documentclass[10pt]{article}
\usepackage{algebra2-article}
\usepackage{algebra2-key}   % pulls in -boxes; do NOT also load -boxes
\newcommand{\TallMath}[1]{$\displaystyle #1\rule[-1.4em]{0pt}{3.2em}$}

\begin{document}

\pageheader{Unit 5, Lesson 5.5}{Exit Ticket --- Answer Key}
\namedateperiod

\vspace{0.06in}

No notes. Factor first --- everything else follows from the factored form.

\vspace{0.04in}

\begin{enumerate}[leftmargin=1.8em, itemsep=8pt]
  \item \textbf{Build it.} \; $f(x)=\dfrac{x^2-1}{x^2+2x-3}$

        \vspace{3pt}
        Factored: $\dfrac{(\text{\ans{$x-1$}})(\text{\ans{$x+1$}})}%
        {(\text{\ans{$x+3$}})(\text{\ans{$x-1$}})}$ \quad Restrictions: $x\neq$ \ans{$1$},
        $x\neq$ \ans{$-3$} \quad Simplified: \ans{$\frac{x+1}{x+3}$}

        \vspace{3pt}
        Hole at \ans{$(1,\frac12)$} \quad Vertical asymptote \ans{$x=-3$} \quad Horizontal asymptote
        \ans{$y=1$}

        \vspace{3pt}
        $x$-intercept \ans{$(-1,0)$} \quad $y$-intercept \ans{$(0,\frac13)$} \quad
        Domain \ans{$(-\infty,-3)\cup(-3,1)\cup(1,\infty)$}

  \item \textbf{Sign chart.} Use the boundary points from item 1 to fill in the table, then answer the
        question below it.
        \begin{center}
        \renewcommand{\arraystretch}{1.4}
        \begin{tabularx}{0.86\linewidth}{|C{2.4cm}|C{1.5cm}|C{1.6cm}|X|}
        \hline
        \rowcolor{forestbg}
        \textbf{Interval} & \textbf{Test $x$} & \textbf{Sign} & \textbf{Above / below} \\ \hline
        $x<-3$    & $-4$ & \ans{$+$} & \ans{above} \\ \hline
        $-3<x<-1$ & $-2$ & \ans{$-$} & \ans{below} \\ \hline
        $x>-1$    & $0$  & \ans{$+$} & \ans{above} \\ \hline
        \end{tabularx}
        \end{center}
        \vspace{-4pt}
        On which interval is the graph \emph{below} the $x$-axis? \ans{$(-3,-1)$} \\[3pt]
        Does this graph ever cross its horizontal asymptote? Answer and explain in one sentence.
        \ansline{No --- $\frac{x+1}{x+3}=1$ would need $x+1=x+3$, which is impossible}

  \item \textbf{SOL-style multiple choice.} Which function has a hole at $x=1$, a vertical asymptote at
        $x=-2$, and a horizontal asymptote $y=1$?
        \begin{enumerate}[label=\Alph*., leftmargin=1.9em, itemsep=1pt]
          \item $y=\dfrac{x-1}{(x-1)(x+2)}$
          \item $y=\dfrac{(x-1)(x+3)}{(x-1)(x+2)}$
                \quad \textcolor{keyred}{\textbf{$\leftarrow$ correct}}
          \item $y=\dfrac{(x+1)(x+3)}{(x+1)(x-2)}$
          \item $y=\dfrac{(x-1)(x+2)}{(x+2)(x+3)}$
        \end{enumerate}
        \vspace{2pt}
        {\small $x-1$ cancels (hole at $1$), $x+2$ survives (wall at $-2$), and the degrees are equal
        with leading coefficients $1$ and $1$ (HA $y=1$). \textbf{A} gets both restrictions right but
        its top degree is lower, so its HA is $y=0$. \textbf{C} has the signs flipped --- hole at $-1$,
        wall at $2$. \textbf{D} cancels the wrong factor: hole at $-2$, wall at $-3$.}
\end{enumerate}

\begin{teachernote}
Graded for completion --- ``mistakes happen, blanks don't.'' Sort into four piles: \textbf{(1) got it};
\textbf{(2) cancel test} --- called $x=1$ a vertical asymptote, or chose \textbf{D}; \textbf{(3) degree
comparison} --- HA given as $y=0$, or chose \textbf{A}; \textbf{(4) sign chart} --- item 2 blank or
signs guessed rather than tested.

\vspace{3pt}
Pile 3 is the fastest fix. \textbf{Pile 2 is the one to re-teach before 5.6}, where the same error --- a
restriction read off the simplified form --- reappears as an extraneous solution and costs a whole
problem. Anyone who wrote $(1,0)$ as an $x$-intercept belongs in pile 2 as well.

\vspace{3pt}
Item 2's last question is the interpretation half, not a bonus. Its answer is \emph{no}, and the reason
matters: once the common factor cancels, this function is a transformed parent again, so the crossing
is impossible for exactly Lesson 5.4's reason. A student who answers ``no, graphs can never cross an
asymptote'' is right for the wrong reason --- correct it out loud, because today's Worked Example A
crossed $y=0$ at its own $x$-intercept.
\end{teachernote}

\end{document}
